\documentclass[a4paper,12pt]{article}
\usepackage[italian]{babel}
\begin{document}
	
	\title{Verifica SQL – Soluzioni}
	\date{}
	\maketitle
	
	\section*{Soluzioni Query}
	
	\begin{enumerate}
		
		\item 
		SELECT nome, cognome, città
		FROM CLIENTI
		WHERE città IN ('Roma','Milano')
		ORDER BY cognome;
		
		\item 
		SELECT AUTO.marca, AUTO.modello, CATEGORIE.costo giornaliero
		FROM AUTO
		JOIN CATEGORIE ON AUTO.id categoria = CATEGORIE.id categoria
		WHERE CATEGORIE.nome categoria = 'SUV';
		
		\item 
		SELECT CLIENTI.nome, CLIENTI.cognome, AUTO.targa, NOLEGGI.data inizio, NOLEGGI.data fine
		FROM NOLEGGI
		JOIN CLIENTI ON NOLEGGI.id cliente = CLIENTI.id cliente
		JOIN AUTO ON NOLEGGI.id auto = AUTO.id auto;
		
		\item 
		SELECT id categoria, COUNT(*) AS numero auto
		FROM AUTO
		GROUP BY id categoria;
		
		\item 
		SELECT DISTINCT CLIENTI.nome, CLIENTI.cognome
		FROM CLIENTI
		JOIN NOLEGGI ON CLIENTI.id cliente = NOLEGGI.id cliente;
		
		\item 
		SELECT id cliente, COUNT(*) AS numero noleggi
		FROM NOLEGGI
		GROUP BY id cliente
		HAVING COUNT(*) > 2;
		
		\item 
		SELECT id noleggio, SUM(importo) AS totale pagato
		FROM PAGAMENTI
		GROUP BY id noleggio;
		
		\item 
		SELECT CLIENTI.id cliente, CLIENTI.nome, CLIENTI.cognome, SUM(PAGAMENTI.importo) AS totale speso
		FROM CLIENTI
		JOIN NOLEGGI ON CLIENTI.id cliente = NOLEGGI.id cliente
		JOIN PAGAMENTI ON NOLEGGI.id noleggio = PAGAMENTI.id noleggio
		GROUP BY CLIENTI.id cliente, CLIENTI.nome, CLIENTI.cognome;
		
		\item 
		SELECT DISTINCT id auto
		FROM MANUTENZIONI;
		
		\item 
		SELECT *
		FROM AUTO
		WHERE id auto NOT IN (
		SELECT id auto
		FROM NOLEGGI
		);
		

		
		\item 
		SELECT id auto, COUNT(*) AS numero noleggi
		FROM NOLEGGI
		GROUP BY id auto;
		
		\item 
		SELECT DISTINCT CLIENTI.nome, CLIENTI.cognome
		FROM CLIENTI
		JOIN NOLEGGI ON CLIENTI.id cliente = NOLEGGI.id cliente
		JOIN AUTO ON NOLEGGI.id auto = AUTO.id auto
		JOIN CATEGORIE ON AUTO.id categoria = CATEGORIE.id categoria
		WHERE CATEGORIE.nome categoria = 'Lusso';
		
		\item 
		SELECT id categoria, AVG(costo giornaliero) AS media costo
		FROM CATEGORIE
		GROUP BY id categoria
		HAVING AVG(costo giornaliero) > 50;
		
		\item 
		SELECT id auto, COUNT(*) AS numero noleggi
		FROM NOLEGGI
		GROUP BY id auto
		HAVING COUNT(*) = (
		SELECT MAX(num)
		FROM (
		SELECT COUNT(*) AS num
		FROM NOLEGGI
		GROUP BY id auto
		) AS Conteggi
		);
		
	\end{enumerate}
	
	\section*{INSERT}
	INSERT INTO CLIENTI (id cliente, nome, cognome, città, patente) \\
	VALUES (21, 'Luca', 'Bianchi', 'Torino', 'AB123CD');
	
	\section*{UPDATE}
	UPDATE CATEGORIE \\
	SET costo giornaliero = costo giornaliero * 1.10 \\
	WHERE nome categoria = 'SUV';
	
	\section*{Vista}
	CREATE VIEW spesa clienti AS
	SELECT C.id cliente, C.nome, C.cognome, SUM(P.importo) AS totale speso
	FROM CLIENTI C
	JOIN NOLEGGI N ON C.id cliente = N.id cliente
	JOIN PAGAMENTI P ON N.id noleggio = P.id noleggio
	GROUP BY C.id cliente, C.nome, C.cognome;
	
\end{document}